\documentclass[../main.tex]{subfiles}
\begin{document}
%==========================================TIÊU ĐỀ TẬP========================================
\begin{table}[h]
    \begin{adjustwidth}{-1cm}{}
        \begin{tabular}{>{\centering\arraybackslash}p{6cm}|>{\raggedright\arraybackslash}p{12cm}}
            \multirow{5}{6cm}
            % Cột bên trái: ÉPISODE và số tập
            {\\[-40pt]
                \humanist\fontsize{20pt}{20pt}\selectfont \centering
                \color{eptype}FOLGE\color{black}\\
                \houseterror\fontsize{72pt}{72pt}\selectfont
                \color{epnum}26\color{black}\\[6pt]
                \cafeteria\fontsize{20pt}{20pt}\selectfont\color{epnum}(M3-6)\color{black}
                \\[18pt]
                % \humanist\fontsize{14pt}{14pt}\selectfont
                % \color{epattr}BIENVENUE, \uline{\color{Banpire} Mel} !
                \sen\fontsize{18pt}{18pt}\selectfont
                % \color{epattr}{Minigeschichten}\\
                % \color{black}
            }
             &
            % Dòng thứ nhất: tên tập tiếng Nhật
            {\fotskip\fontsize{18pt}{18pt}\selectfont \ruby{起}{お}こさないようにそーっと！}                                                                                                                                \\[6pt]
            % Dòng thứ hai: tên tập tiếng Anh
             & {\cafeteria\fontsize{18pt}{18pt}\selectfont Sleepy Time}                                                                                                                                       \\[4pt]
            % Dòng thứ ba: tên tập tiếng Việt
             & {\vnmsans\fontsize{18pt}{18pt}\selectfont Bài thơ lục bát về những anh chàng ngủ gật}                                                                                                          \\[8pt]
            % Dòng thứ tư: Nhân vật xuất hiện
             & {\caviar{\fontsize{14pt}{14pt}\selectfont Nhân vật: \emp{kaiho2} \emp{naomi2} \emp{miyabi} \emp{izuru} \emp{aruran1} \emp{rikka} \emp{astel} \emp{temma} \emp{roberu} \emp{shien} \emp{oga}}} \\
            % Dòng thứ năm (nếu có): Nhân vật chỉ xuất hiện trong Omake
             & {\abv Truyện phụ:} \emp{kuon2} \emp{ryuuhou2} \emp{suzuno2}
            \\[8pt]
            % Dòng cuối cùng: Thời gian diễn ra của tập (thường sẽ là ngày phát hành)
             & {\sen\fontsize{16pt}{16pt}\selectfont Mốc thời gian: 22/07/2022}                                                                                                                               \\[18pt]
        \end{tabular}\\[8pt]
    \end{adjustwidth}
\end{table}
%=========================================TRUYỆN CHÍNH========================================
\color{black}

\txtbx[2]{
    {\color{vol} \uline{Tóm tắt:}} Những mẩu chuyện dở khóc dở cười xoay quanh các chàng trai \hls\ khi ngủ gật cùng sự chứng kiến của hai chị em Kaiho và cô bé Naomi mang theo chiếc máy ảnh nhỏ: Izuru ngủ gục trên bàn khiến Aruran muốn ra dáng đàn anh nhưng lại túm tóc cậu nên liên tiếp nhận hai cú đấm trời giáng, trong khi Naomi ví von Izuru ngoan ngoãn như một chú mèo con nếu được xoa đầu đúng cách; Temma ngủ sấp chảy dãi đè lên chiếc điện thoại của Roberu, Astel tốt bụng nhấc bổng đầu Temma như bao khoai tây để lấy máy hộ nhưng lại bị Roberu quất quạt tới tấp vì cách đối xử thô bạo với bạn; Oga ngủ say trên bàn ăn bị Shien táy máy làm gãy mất chiếc sừng, sau một phen hú vía vái lạy gá sừng bất thành, Shien ngỡ sẽ bị ăn đòn lôi đình nhưng hóa ra sừng của Quỷ tướng có thể mọc lại cái mới khi được tưới đủ nước và thủ thỉ những lời dịu dàng, khiến Kaiho ngán ngẩm bình phẩm chẳng khác nào việc chăm cây; và việc mà các chàng trai phản ứng khi Kaiho ngủ say vì mệt mỏi, để rồi bị cô em út tinh nghịch chụp trộm gửi về nhóm gia đình.
}

\begin{center}
    \uline{(Như tiêu đề của tập, câu chuyện này sẽ được viết dưới dạng thể thơ lục bát, vì tập này không có thoại, ngoại trừ của Kaiho và Naomi.)}
\end{center}

\ct

\begin{center}
    Bàn con góc nhỏ dịu êm,\\
    Nhất Huyền\footnote{Tức Izuru, lấy từ Hán-Việt 一弦.} gục xuống ngủ thêm giấc nồng.\\
    Đôi tay làm gối thong dong,\\
    Hải Phàm\footnote{Tức Kaiho, lấy từ Hán-Việt 海帆.} ngồi cạnh chăm lòng sổ biên.\\
    Trực Mỹ\footnote{Tức Naomi, lấy từ Hán-Việt 直美.} máy ảnh mang bên,\\
    Nụ cười trong trẻo an nhiên sắc màu.\\
    Bỗng đâu A Nhĩ\footnote{Tức Aruran, lấy từ A Nhĩ Lan 阿爾蘭.} bước mau,\\
    Nở nụ cười ấm ngọt ngào mẹ cha.\\
    Tiến gần đưa ngón tay ra,\\
    Nắm ngay chỏm tóc ngỡ là vuốt chơi.\\
    Nhất Huyền mở mắt tức thời,\\
    Vung tay giáng thẳng một quyền cực nhanh!\\
    A-rư xoay tít vòng quanh,\\
    Đổ nhào xuống đất tan tành tấm thân.\\
    Chàng Việc Vặt lại vùi chân,\\
    Chìm sâu giấc ngủ chẳng cần bận tâm.\\
    Hải Phàm ngán ngẩm buông lời,\\
    Trách người dại dột chọc người lúc mê.\\
    Trực Mỹ cười khẽ tỉ tê:\\
    ``Anh ấy như một chú mèo nhỏ thôi!\\
    Mèo thì vuốt nhẹ êm trôi,\\
    Ai đi túm tóc giật chồi thế kia!''\\
    A-rư bất tỉnh nằm lì,\\
    Cả gian nín lặng tứ bề xót xa.

    \begin{figure}[H]
        \centering
        \begin{subfigure}[t]{0.45\textwidth}
            \centering
            \includegraphics[width=8cm]{../../../../Image/BookM3/m3-6a1.png}
        \end{subfigure}
        \quad
        \begin{subfigure}[t]{0.45\textwidth}
            \centering
            \animategraphics[autoplay,loop,width=8cm]{12}{../../../../Image/BookM3/m3-6a2/m3-6a2.}{1}{89}
        \end{subfigure}
    \end{figure}

    Luật Khả\footnote{Tức Rikka, lấy từ Hán-Việt 律可.}, Nhã\footnote{Tức Miyabi, lấy từ Hán-Việt 雅.} đứng bên nhà,\\
    Mồ hôi nhỏ giọt thật thà sợ lây.\\
    Hải Phàm khoát nhẹ bàn tay,\\
    Liếc nhìn cảnh báo chớ gây họa tai.\\
    Nhã bèn rón rén một hai,\\
    Đưa tay khẽ vuốt tóc dài người ta.\\
    Lạ thay chẳng chút phong ba,\\
    Nhất Huyền để mặc Nhã xoa mái đầu.\\
    Cậu chàng Việc Vặt ngủ sâu,\\
    Ngoan như mèo nhỏ ngả đầu bình yên.\\
    Nhã vui hớn hở cười duyên,\\
    Khả nhìn mừng rỡ tính liền học theo.\\
    Vừa đưa tay định xoa ngay,\\
    Nhất Huyền động nhẹ khiến chàng rụt tay.\\
    Chờ khi cậu ngủ mê say,\\
    Khả bèn vuốt nhẹ đôi tay lên đầu.\\
    Thấy người ngủ thiếp bền lâu,\\
    Cả hai nhẹ nhõm gật đầu mừng vui.

    \begin{figure}[H]
        \centering
        \begin{subfigure}[t]{0.45\textwidth}
            \centering
            \includegraphics[width=8cm]{../../../../Image/BookM3/m3-6b1.png}
        \end{subfigure}
        \quad
        \begin{subfigure}[t]{0.45\textwidth}
            \centering
            \includegraphics[width=8cm]{../../../../Image/BookM3/m3-6b2.png}
        \end{subfigure}
        \quad
        \begin{subfigure}[t]{0.45\textwidth}
            \centering
            \includegraphics[width=8cm]{../../../../Image/BookM3/m3-6b3.png}
        \end{subfigure}
    \end{figure}

    Trực Mỹ ánh mắt rạng ngời,\\
    Cũng mon men lại mỉm cười bước sang.\\
    Mặc cô chị cản vội vàng,\\
    Đưa tay vuốt nhẹ dịu dàng tóc êm.\\
    Lại giương máy ảnh chụp liền,\\
    Nhất Huyền khoan khoái ngủ yên gối đầu.\\
    Thấy lòng nhẹ nhõm thảnh thơi,\\
    Cả ba đứng ngắm một hồi an nhiên.\\
    A-rư lồm cồm bò lên,\\
    Thấy ba người vuốt dịu hiền chẳng sao.\\
    Tính làm một cú ngọt ngào,\\
    Nhưng tay quen thói túm vào tóc ngay!\\
    Nào ngờ vừa nắm ngón tay,\\
    Cú quyền giáng thẳng lăn quay ra sàn!\\
    Khả, Nhã trố mắt kinh hoàng,\\
    Hải Phàm lắc nhẹ thở than buông lời:\\
    ``Đã bảo tay ráp giấy bồi,\\
    Lại còn túm tóc, ăn hồi đòn cay!''\\
    Trực Mỹ cười nép sau tay:\\
    ``Chọc mèo nổi đóa, ráng hầu cước bay\dots''

    \gif{m3-6c}{1}{57}

\end{center}

\ct

\begin{center}
    Bàn ăn trưa ngả bóng tà,\\
    Thiên Chân\footnote{Tức Temma, lấy từ Hán-Việt 天真.} nằm sấp như là khúc cây.\\
    Mắt he hé, miệng dãi đầy,\\
    Tướng nằm xấu xí phơi bày cười đau.\\
    Bối Nhĩ\footnote{Tức Roberu, lấy từ Hán-Việt (羅)貝爾.}, Trực Mỹ nhìn nhau,\\
    Bụm môi nín thở cúi đầu cười rung.\\
    Hải Phàm cất tiếng nhắc chừng:\\
    ``Hiệp sĩ mà ngủ lạ lùng thế kia!''\\
    Bối Nhĩ toan tính chụp ngay,\\
    Lục tung túi áo tìm tay máy mình.\\
    A Tư\footnote{Tức Astel, lấy từ Hán-Việt 阿斯(特爾).} chợt bước lặng thinh,\\
    Vỗ vai Bối Nhĩ chỉ hình dưới thân.\\
    Lưng cam chiếc máy rơi gần,\\
    Ngay dưới trán bạn đang phần đè lên.\\
    Bối Nhĩ cười gượng ngó lên,\\
    A Tư đã bước lại bên tức thì.

    \begin{figure}[H]
        \centering
        \begin{subfigure}[t]{0.45\textwidth}
            \centering
            \includegraphics[width=8cm]{../../../../Image/BookM3/m3-6d1.png}
        \end{subfigure}
        \quad
        \begin{subfigure}[t]{0.45\textwidth}
            \centering
            \includegraphics[width=8cm]{../../../../Image/BookM3/m3-6d2.png}
        \end{subfigure}
    \end{figure}

    Ba người nín thở sợ nguy,\\
    Xem chàng Bảnh Bao tính gì ra tay.\\
    Cậu chàng xốc nhẹ đôi tay,\\
    Nhấc bổng đầu bạn tựa bì khoai tây!\\
    Nửa say nửa tỉnh ngất ngây,\\
    Miệng co giật méo dãi đầy hai bên.\\
    Hải Phàm bật tiếng cười nghiêng,\\
    A Tư rút máy đặt liền xuống ngay.\\
    Đặt đầu cái bịch buông tay,\\
    Thiên Chân vẫn ngủ mê say giấc nồng.\\
    A Tư hớn hở trong lòng,\\
    Đem trao lại máy tươi hồng mặt hoa.\\
    Nào ngờ Bối Nhĩ rút ra,\\
    Cây quạt giấy quất thật thà giáng ngay!\\
    Quát rằng: ``Nhấc tựa cỏ cây,\\
    Xem đầu bạn hữu như bầy khoai thô!?''

    \begin{figure}[H]
        \centering
        \begin{subfigure}[t]{0.45\textwidth}
            \centering
            \includegraphics[width=8cm]{../../../../Image/BookM3/m3-6e1.png}
        \end{subfigure}
        \quad
        \begin{subfigure}[t]{0.45\textwidth}
            \centering
            \animategraphics[autoplay,loop,width=8cm]{12}{../../../../Image/BookM3/m3-6e2/m3-6e2.}{1}{57}
        \end{subfigure}
    \end{figure}

    Cú quạt giáng xuống bất ngờ,\\
    A Tư ngã ngửa bơ phờ ngã lăn!\\
    Chiếc máy tung bổng xoay vần,\\
    May sao Bối Nhĩ nhanh chân chộp liền.\\
    Nhìn sang tướng ngủ bạn hiền,\\
    Dù lòng phát tởm mặt liền đanh trơ.\\
    Bối Nhĩ chẳng bỏ lỡ thời,\\
    Bấm ngay một phát để đời ảnh dìm!\\
    A Tư ngồi dậy im lìm,\\
    Xoa đầu ngơ ngác mắt chìm hỏi han:\\
    Lập công sao chịu tiếng oan?\\
    Hải Phàm nhếch mép: ``Đáng vì thô hung!''\\
    Trực Mỹ cười khẽ ung dung:\\
    ``Giúp người kiểu đó, ăn cùng đòn đau\dots''

    \img{m3-6f}

\end{center}

\ct

\begin{center}
    Bàn ăn tĩnh lặng chiều vàng,\\
    Hoàng Nha\footnote{Tức Oga, lấy từ Hán-Việt 荒咬.} nằm ngủ mơ màng say sưa.\\
    Tử Diễm\footnote{Tức Shien, lấy từ Hán-Việt 影山.} ngồi ngắm đong đưa,\\
    Thò tay mân mó sừng vừa bạn ta.\\
    Bé Trực khẽ nhắc nhỏ qua:\\
    ``Chớ sờ kẻo thức người ta dậy liền!''\\
    Diễm xua tay trấn an yên:\\
    ``Em sờ nhè nhẹ làm phiền chi đâu.''\\
    Hải Phàm cảnh báo từ đầu:\\
    ``Coi chừng rước họa chuốc sầu vào thân!''

    \img{m3-6g}

    Bàn tay táy máy nhiều lần,\\
    Bỗng nghe cái rắc, sừng dần rơi ngay!\\
    Sừng rơi lăn lóc khỏi mày,\\
    Tử Diễm mặt tái trắng bày như vôi.\\
    Bé Trực toát lạnh mồ hôi,\\
    Sợ người tỉnh giấc cuốn trôi gian phòng.\\
    Hải Phàm ngán ngẩm trong lòng:\\
    ``Chị đây đã dặn coi chừng chẳng nghe!\\
    Nó mà mở mắt nhìn đe,\\
    Là tan tành xác chớ khoe mạng này!''\\
    Diễm run lẩy bẩy hai tay,\\
    Vội vàng nhặt lấy chắp ngay lên đầu.\\
    Rón rén từng bước lui sau,\\
    Hai tay hướng phía sừng mau niệm cầu:\\
    Xin sừng yên vị bền lâu,\\
    Đừng rơi xuống nữa thêm sầu dạ tôi!

    \begin{figure}[H]
        \centering
        \begin{subfigure}[t]{0.45\textwidth}
            \centering
            \includegraphics[width=8cm]{../../../../Image/BookM3/m3-6h1.png}
        \end{subfigure}
        \quad
        \begin{subfigure}[t]{0.45\textwidth}
            \centering
            \includegraphics[width=8cm]{../../../../Image/BookM3/m3-6h2.png}
        \end{subfigure}
    \end{figure}

    Thấy sừng nằm lại êm trôi,\\
    Cậu thở phào nhẹ, nụ cười dần tươi.\\
    Đôi tai thú cụp thảnh thơi,\\
    Bé Trực vội bước lại nơi chị mình.\\
    Bỗng nghe lách cách rung rinh,\\
    Sừng rơi rớt lại thình lình xuống sàn!\\
    Chưa kịp đón lấy hay can,\\
    Thì Hoàng Nha đã mơ màng nhìn sang.\\
    Thấy Diễm quỳ mọp hai hàng,\\
    Hai tay nâng chiếc sừng vàng van xin.\\
    Nét mặt tội nghiệp lặng thinh,\\
    Đôi tai thú cụp van xin tha mình.\\
    Tưởng đâu nổi trận lôi đình,\\
    Quát vang giáng xuống tội tình đứa em.\\
    Nào ngờ Hoàng Nha nhìn xem,\\
    Thở dài một tiếng êm đềm bảo ban:\\
    ``Gãy rồi chớ có thở than,\\
    Sừng này mọc lại dễ dàng có sao!\\
    Mau mau rót chút nước vào,\\
    Thêm lời âu yếm ngọt ngào bên tai!''

    \img{m3-6i}

    Diễm nghe mắt sáng hân hoan,\\
    Cùng Trực vội rót nước tràn gốc chân.\\
    Thủ thỉ lời ngọt ân cần:\\
    ``Sừng ngoan mau nhú, mau phần lớn nhanh!''\\
    Kỳ thay từ chỗ gãy lành,\\
    Chiếc sừng vươn vút mọc nhanh tức thời!\\
    Diễm mừng thở nhẹ thảnh thơi,\\
    Bé Trực bấm máy lưu thời khắc hay.\\
    Hải Phàm khoanh nhẹ đôi tay,\\
    Lắc đầu bình phẩm chuyện hay chốn này:\\
    ``Chẳng khác nào việc chăm cây,\\
    Tưới nước thủ thỉ, sừng đây mọc liền!''
\end{center}

%==========================================TRUYỆN PHỤ=========================================
\omk

\begin{center}
    Chiều tàn buông ánh hoàng hôn,\\
    Hải Phàm mệt lả vùi chôn giấc nồng.\\
    Cả ngày sổ sách bận lòng,\\
    Răn đe lũ nghịch tốn công nhọc nhằn.\\
    Gục đầu xuống mép bàn ăn,\\
    Nàng buông giấc ngủ chẳng cần bận tâm.

    Lũ trai ngơ ngác nhìn trân,\\
    Bà chằn Tổng quản đang phần ngủ say!\\
    Xưa nay thép lệnh trong tay,\\
    Thét ra lửa cháy chẳng ai dám nhờn.\\
    Bối Nhĩ toan tính bước sang,\\
    Định thò ngón chọc trêu nàng một phen.\\
    Nào ngờ vừa nhích chân chen,\\
    Hải Phàm chợt nhíu đôi viền mi cong.\\
    Sát khí tỏa lạnh trong lòng,\\
    Bàn tay nắm chặt lấy vòng thước gai!\\
    Lũ trai hồn vía lạc loài,\\
    Vội vàng thụt lại né dài ra xa.\\
    Thầm răn: ``Chớ đụng người ta,\\
    Đụng vào cọp cái khéo là mạng vong!''

    Bé Trực khẽ bước thong dong,\\
    Lấy manh áo khoác ấm lòng đắp lên.\\
    Hơi ấm lan tỏa dịu hiền,\\
    Chân mày giãn nhẹ, ngủ yên giấc nồng.\\
    Nét mặt trở lại trắng hồng,\\
    Khác xa vẻ dữ gắt gồng ngày thường.\\
    Bé con mỉm nụ cười thương,\\
    Nâng niu máy ảnh bấm liền góc xinh.\\
    Tách! một tiếng trúng khung hình,\\
    Gửi về nhóm nhỏ gia đình xem chung.

    Chuông reo phá vỡ khoảng không,\\
    Cả nhà rôm rả ấm lòng trêu vui.\\
    Lưu Bảo bật tiếng cười tươi:\\
    ``Bà chằn cũng có phút người nết na!''\\
    Linh Nãi\footnote{Tức Suzuno, lấy từ Linh Nãi 鈴乃.} khẽ thốt lời hoa:\\
    ``Thương chị Hai quá, mệt nhoài giấc say\dots''\\
    Cửu Viễn ngắm ảnh cười say,\\
    Nhớ hồi Sấm Âm ngủ ngày trong phòng\footnote{{\flaregothic ÉPISODE 126}, quyển số Mười. Về các tên Lưu Bảo và Cửu Viễn, xem lại tập đó.}:\\
    ``May cho lũ trẻ biết lòng,\\
    Chứ hùa nghịch ngợm lòng vòng chọc chơi.\\
    Như hồi trêu Chó khắp nơi,\\
    Xoa đầu, đút bánh, trêu người ngủ say\dots\\
    Gặp Hải Phàm tỉnh phen này,\\
    Ắt là ăn cước lăn quay cả đàn!''

    Tiếng chuông tin nhắn rộn ràng,\\
    Khiến cho nàng Tổng mơ màng tỉnh ra.\\
    Mở xem bức ảnh gửi qua,\\
    Kèm tràng bình luận cả nhà trêu vui.\\
    Mặt nàng đỏ ửng một hồi,\\
    Chẳng hề quát mắng buông lời dọa đe.\\
    Chỉ ôm lấy mặt thẹn che,\\
    Nhìn cô em út khẽ hoe ngượng ngùng:\\
    ``Bé con nghịch ngợm quá chừng,\\
    Lần sau chớ chụp chớ mừng gửi khoe!''\\
    Bé Trực khúc khích cười xòe,\\
    Nép vào lòng chị ấm che ân tình.\\
    Dẫu ngoài nghiêm nghị một mình,\\
    Vẫn vui để thấy êm đềm tình thân.\\
    Ở nhà Cửu Viễn thấu tâm,\\
    Nhìn hình em gái mỉm thầm nụ hoa\dots
\end{center}

\end{document}